\section{Diseño de base de datos}
\label{sec:basededatos}

\subsection{Modelo entidad-relación}

\begin{figure}[htbp]
\centering
\resizebox{\textwidth}{!}{% Modelo entidad-relación de PayCore (simplificado a lo esencial).
\begin{tikzpicture}[
  font=\scriptsize,
  >={Stealth[length=2mm]},
  entidad/.style={draw=acento!55, fill=acento!7, rounded corners=3pt,
                  text width=40mm, inner sep=5pt},
  titulo/.style={font=\scriptsize\bfseries, color=acento},
  rel/.style={->, draw=tinta!55},
  nota/.style={font=\scriptsize\itshape, color=tinta!65},
]

\node[entidad] (users) {%
  \textcolor{acento}{\bfseries users}\\[2pt]
  \codigo{id} \textsc{uuid pk}\\
  \codigo{email} \textsc{citext uk}\\
  \codigo{password\_hash}\\
  \codigo{full\_name}\\
  \codigo{created\_at}};

\node[entidad, right=26mm of users] (accounts) {%
  \textcolor{acento}{\bfseries accounts}\\[2pt]
  \codigo{id} \textsc{uuid pk}\\
  \codigo{user\_id} \textsc{fk uk}\\
  \codigo{alias} \textsc{uk}\\
  \codigo{kind} user\,|\,system\\
  \codigo{currency} = PEN\\
  \codigo{balance\_cents} \textsc{bigint}\\
  \codigo{created\_at}};

\node[entidad, below=14mm of users] (idem) {%
  \textcolor{acento}{\bfseries idempotency\_keys}\\[2pt]
  \codigo{id} \textsc{uuid pk}\\
  \codigo{user\_id} \textsc{fk}\\
  \codigo{key}\\
  \codigo{request\_fingerprint}\\
  \codigo{state}\\
  \codigo{response\_status}\\
  \codigo{response\_body} \textsc{jsonb}\\
  \codigo{transfer\_id} \textsc{fk}\\
  \codigo{expires\_at}};

\node[entidad, below=14mm of accounts] (transfers) {%
  \textcolor{acento}{\bfseries transfers}\\[2pt]
  \codigo{id} \textsc{uuid pk}\\
  \codigo{source\_account\_id} \textsc{fk}\\
  \codigo{destination\_account\_id} \textsc{fk}\\
  \codigo{amount\_cents} \textsc{bigint}\\
  \codigo{status}\\
  \codigo{rejection\_reason}\\
  \codigo{concept}\\
  \codigo{completed\_at}};

\node[entidad, right=22mm of transfers] (ledger) {%
  \textcolor{acento}{\bfseries ledger\_entries}\\[2pt]
  \codigo{id} \textsc{uuid pk}\\
  \codigo{transfer\_id} \textsc{fk}\\
  \codigo{account\_id} \textsc{fk}\\
  \codigo{amount\_cents} \textsc{bigint}\\
  \hspace{4pt}{\itshape signo: $-$ débito, $+$ crédito}\\
  \codigo{balance\_after\_cents}\\
  \codigo{created\_at}};

\draw[rel] (users.east) -- node[above] {1\,:\,1} (accounts.west);
\draw[rel] (users.south) -- node[right] {1\,:\,N} (idem.north);
\draw[rel] (accounts.south) -- node[right] {1\,:\,N} (transfers.north);
\draw[rel] (transfers.east) -- node[above, align=center]
  {1\,:\,2\\[-2pt]{\itshape exactamente}} (ledger.west);
\draw[rel] (idem.east) -- node[above] {0\,:\,1} (transfers.west);

\node[nota, below=5mm of transfers.south, text width=140mm,
      align=left, anchor=north] {%
  El libro mayor (\codigo{ledger\_entries}) es de solo inserción y es la
  fuente de verdad. La columna \codigo{accounts.balance\_cents} es una
  proyección materializada que se escribe en la misma transacción. La
  restricción de no negatividad exime a las cuentas de sistema, que son el
  origen auditado del saldo ficticio.};

\end{tikzpicture}
}
\caption{Modelo de datos de PayCore.}
\label{fig:er}
\end{figure}

\subsection{El libro mayor como fuente de verdad}

Cada transferencia completada genera \textbf{exactamente dos asientos} en
\codigo{ledger\_entries}: uno negativo en la cuenta de origen y uno positivo
en la de destino, de igual magnitud y ligados por \codigo{transfer\_id}. La
tabla es de solo inserción.

La columna \codigo{accounts.balance\_cents} es una \emph{proyección
materializada} de ese libro mayor, escrita en la misma transacción. No es una
duplicación descuidada sino una decisión con contrapartidas explícitas:

\begin{itemize}[leftmargin=1.2em]
  \item \textbf{A favor:} lectura de saldo en $O(1)$, existe una fila concreta
    sobre la que adquirir el bloqueo, y la restricción
    \codigo{CHECK (balance\_cents >= 0)} es posible porque hay una columna que
    restringir.
  \item \textbf{En contra:} dos representaciones que deben coincidir. Se
    acepta con dos obligaciones: escribir ambas en la misma transacción, y una
    prueba de reconciliación que recalcula desde el libro mayor y se ejecuta
    al final de cada suite y después de cada carrera de carga.
\end{itemize}

\subsection{El origen del dinero simulado}

El saldo ficticio no aparece de la nada. Entra desde una cuenta de sistema,
\codigo{SYSTEM\_FUNDING}, que es la \textbf{única exenta} de la restricción de
no negatividad. Así la suma global de asientos sigue siendo exactamente cero
incluso después de cargar los datos de prueba, y el origen queda auditado en
lugar de materializarse por decreto. Es la excepción declarada a la
conservación del saldo total, y es una operación administrativa explícita que
no se expone por la API. \estado{verificado}

\subsection{Invariantes y dónde se aplican}

El cuadro~\ref{tab:invariantes} enumera las diez invariantes del sistema. La
columna más importante es la tercera: nueve de las diez se aplican en el motor
de base de datos, no en Python.

\begin{table}[htbp]
\centering
\caption{Invariantes del sistema y su punto de aplicación.}
\label{tab:invariantes}
\scriptsize
\begin{tabularx}{\textwidth}{@{}llXl@{}}
\toprule
\textbf{Id} & \textbf{Invariante} & \textbf{Dónde se aplica} &
\textbf{Prueba} \\
\midrule
I1 & Ningún saldo de usuario negativo &
  \codigo{CHECK (kind = 'system' OR balance\_cents >= 0)} & B y D \\
I2 & Importe estrictamente positivo &
  \codigo{CHECK (amount\_cents > 0)} y validación de esquema & F \\
I3 & Origen distinto de destino & \codigo{CHECK} de tabla & unitaria \\
I4 & Dos asientos por transferencia que suman cero &
  Trigger de restricción \codigo{DEFERRABLE INITIALLY DEFERRED} & A y esquema \\
I5 & Suma global del libro mayor igual a cero &
  Prueba de reconciliación sobre la base & toda la suite \\
I6 & Saldo materializado igual a la suma de asientos &
  Prueba de reconciliación & toda la suite \\
I7 & Ninguna transferencia parcial & Una sola transacción SQL & E \\
I8 & Un reintento no crea movimientos &
  Clave persistida en la misma transacción & C y E \\
I9 & Un usuario solo opera sobre su cuenta &
  La cuenta origen sale del token, nunca del cuerpo & F \\
I10 & El libro mayor es de solo inserción &
  Trigger que rechaza \codigo{UPDATE} y \codigo{DELETE} & esquema \\
\bottomrule
\end{tabularx}
\end{table}

\subsection{Dos garantías que merecen detalle}

\paragraph{I4, verificada al confirmar.} Un \codigo{CHECK} de columna no puede
expresar \enquote{esta transferencia tiene exactamente dos asientos que suman
cero}, porque es una condición sobre un conjunto de filas que se insertan
dentro de la misma transacción. Se implementa con un trigger de restricción
diferido: dentro de la transacción los asientos pueden estar momentáneamente
descuadrados, y es el \codigo{COMMIT} el que rechaza el conjunto si no cuadra.

\begin{lstlisting}[language=SQL,caption={Trigger de restricción que verifica
el cuadre al confirmar la transacción.},label={lst:trigger}]
CREATE OR REPLACE FUNCTION paycore_check_ledger_balanced() RETURNS trigger AS $$
DECLARE v_count integer; v_sum bigint;
BEGIN
    SELECT count(*), coalesce(sum(amount_cents), 0) INTO v_count, v_sum
      FROM ledger_entries WHERE transfer_id = NEW.transfer_id;
    IF v_count <> 2 THEN
        RAISE EXCEPTION 'ledger desbalanceado: la transferencia % tiene % asientos
            y se esperaban 2', NEW.transfer_id, v_count
            USING ERRCODE = 'check_violation';
    END IF;
    IF v_sum <> 0 THEN
        RAISE EXCEPTION 'ledger desbalanceado: los asientos de la transferencia %
            suman % y se esperaba 0', NEW.transfer_id, v_sum
            USING ERRCODE = 'check_violation';
    END IF;
    RETURN NULL;
END;
$$ LANGUAGE plpgsql;

CREATE CONSTRAINT TRIGGER trg_ledger_balanced
AFTER INSERT ON ledger_entries
DEFERRABLE INITIALLY DEFERRED
FOR EACH ROW EXECUTE FUNCTION paycore_check_ledger_balanced();
\end{lstlisting}

\paragraph{I10, solo inserción.} Un segundo trigger rechaza cualquier
\codigo{UPDATE} o \codigo{DELETE} sobre \codigo{ledger\_entries}. La operación
\codigo{TRUNCATE} sigue siendo posible, que es lo que usan las pruebas para
limpiar entre casos; una corrección contable se haría con un asiento de
reversión, nunca borrando historia.

\subsection{Migraciones}

El esquema completo se crea en una única migración de Alembic versionada,
incluidos los dos triggers y la extensión \codigo{citext}. Dos pruebas
protegen esta parte: una comprueba que \codigo{downgrade base} seguido de
\codigo{upgrade head} reconstruye el esquema, y otra que el esquema migrado no
diverge de los modelos declarados, de modo que cambiar un modelo y olvidar la
migración pone el CI en rojo. \estado{verificado}
